% !TEX root = ../main.tex
% =========================================================
% daftarsingkatan.tex - DAFTAR SINGKATAN DAN NOTASI
% =========================================================

\clearpage
\phantomsection
\addcontentsline{toc}{chapter}{DAFTAR SINGKATAN}

\begin{center}
    \large \textbf{DAFTAR SINGKATAN DAN NOTASI}
\end{center}
\vspace{3em}

\begin{center}
    \begin{tabularx}{0.8\textwidth} {
            >{\raggedright\arraybackslash}X
            >{\raggedright\arraybackslash}X}
        \toprule
        \textbf{Notasi} & \textbf{Arti} \\
        \midrule
        PV       & \textit{Photovoltaic} \\
        PLTS     & Pembangkit Listrik Tenaga Surya \\
        STC      & \textit{Standard Test Conditions} \\
        OPC      & \textit{Open-Circuit} \\
        I--V     & Arus--Tegangan (\textit{current--voltage}) \\
        P--V     & Daya--Tegangan (\textit{power--voltage}) \\
        MPP      & \textit{Maximum Power Point} \\
        MPPT     & \textit{Maximum Power Point Tracking} \\
        O\&M     & \textit{Operation and Maintenance} \\
        IEC      & \textit{International Electrotechnical Commission} \\
        $V_{oc}$ & Tegangan \textit{open-circuit} \\
        $I_{sc}$ & Arus \textit{short-circuit} \\
        $V_{mp}$ & Tegangan pada titik daya maksimum \\
        $I_{mp}$ & Arus pada titik daya maksimum \\
        $P_{max}$& Daya maksimum modul \\
        $FF$     & \textit{Fill Factor} \\
        $R_{VM}$ & Rasio $V_{mp}/V_{oc}$ (indikator sisi tegangan) \\
        $R_{IM}$ & Rasio $I_{mp}/I_{sc}$ (indikator sisi arus) \\
        $R_{s}$  & Resistansi seri \\
        $R_{sh}$ & Resistansi shunt \\
        \bottomrule
    \end{tabularx}
\end{center}